% Deokhyun Ahn's downloadable CV.
% Layout supplied by the user in Resume.zip (cv_4.tex + resume.cls).
% Keep resume.cls beside this file when uploading to Overleaf.
% Build: node bin/export_cv.cjs assets/pdf/cv.pdf
% The website's /cv/ page continues to use _data/cv.yml.

\documentclass{resume}
\usepackage{enumitem}
\hypersetup{
  pdftitle={Deokhyun Ahn | Curriculum Vitae},
  pdfauthor={Deokhyun Ahn},
  pdfsubject={Research in reliable and efficient video understanding for social good},
  pdfkeywords={Video understanding, Video anomaly detection, Vision-language models}
}
\newcommand{\cvParagraph}[1]{#1\par\smallskip}
\newcommand{\publication}[3]{%
  \item\begin{minipage}[t]{\linewidth}
    #2 #1 \textit{#3.}
  \end{minipage}
}

\name{Deokhyun Ahn}
\address{+82 10 3733 3254 \\
  \href{mailto:dhan@kist.re.kr}{\nolinkurl{dhan@kist.re.kr}} \\
  \href{mailto:dhahn@korea.ac.kr}{\nolinkurl{dhahn@korea.ac.kr}}}
\address{\href{https://deo-ahn.github.io/}{deo-ahn.github.io} \\
  \href{https://www.linkedin.com/in/deokhyun-ahn/}{LinkedIn} \\
  \href{https://scholar.google.com/citations?user=5VDH7PgAAAAJ&hl=ko}{Google Scholar} \\
  Seoul, Republic of Korea}
\address{Ph.D. Candidate, Brain \& Cognitive Engineering}

\begin{document}
{\centering\small Expected graduation: August 2027
\textbar\ Available for employment from March 2027\par}

\begin{rSection}{Research Interests}
\cvParagraph{\textbf{\textit{Reliable and efficient video understanding for social good.}}}
\cvParagraph{Research for social good has always been my first priority, and it will continue to guide my work.
I am a Ph.D. candidate in \href{https://bce.korea.ac.kr/bce/index.do}{Brain \& Cognitive Engineering at Korea University}
and a student researcher at \href{https://www.kist.re.kr/}{KIST AI\textperiodcentered Robotics}.}
\cvParagraph{My early work on sensing applications, including an ultrasonic SLAM helmet for firefighters and
indoor drone tracking, shaped my commitment to dependable systems for real-world needs. I now study reliable
and efficient video understanding: identifying anomalous behavior, understanding its temporal context, and
explaining it with visual evidence.}
\cvParagraph{My work connects video anomaly detection and understanding, adaptive video sampling, action
recognition, and vision-language models (VLMs). Through public-safety research at KIST, I work on contextual
anomaly detection and behavior analysis for real-world surveillance video.}
\cvParagraph{Recent work includes \href{https://deo-ahn.github.io/buster/}{BUSTER} for VLM-guided anomaly
detection with adaptive sampling, \href{https://ieeexplore.ieee.org/document/11217333}{Where and What} for contextual
anomaly detection, and \href{https://skill-lens.net/}{Look Closer} for action-guided proficiency estimation.}
\cvParagraph{\textbf{Current focus:} why vision-language models hallucinate in video anomaly understanding,
and how selective visual verification can make their explanations more faithful.}
\cvParagraph{\textbf{Research keywords:} Video understanding; Video anomaly detection \& understanding;
Adaptive video sampling; Efficient vision-language models; VLM hallucination \& reliability;
Temporal \& contextual reasoning; Action recognition \& human motion analysis;
Video surveillance \& public safety.}

\end{rSection}

\begin{rSection}{Education}
\begin{rDated}{Mar. 2023 -- \\ Present}{Korea University}{Seoul, KR}{Ph.D. Candidate, Brain \& Cognitive Engineering}
  \item Expected graduation: August 2027.
  \item Study contextual video understanding inspired by human cognitive processes.
\end{rDated}
\begin{rDated}{Feb. 2020 \\ Aug. 2015}{Kwangwoon University}{Seoul, KR}{M.Eng. in Electronic Engineering, February 2020 \\ B.Eng. in Electronic Engineering, August 2015}
  \item Thesis: \emph{A Study on Passive Target Tracking in Three-Dimensional Indoor Environment by Using Deep Neural Network}.
  \item Completed 19 months of mandatory military service.
\end{rDated}
\end{rSection}

\newpage
\begin{rSection}{Research \& Work Experience}
\begin{rDated}{Jan. 2023 -- \\ Present}{Student Researcher}{Seoul, KR}{KIST AI\textperiodcentered Robotics Institute \textbar\ Vision Intelligence Group}
  \item Video anomaly understanding; Surveillance video; Vision language models.
  \item Develop public safety video systems spanning action recognition, contextual anomaly detection,
  and vision language models that explain anomalous behavior. [1] [2] [3] [4] [5]
  \item Developed research prototypes and integrated them into a real-world monitoring system.
  \item Served as frontend team lead for the monitoring system, contributing directly to implementation
  and coordinating the team's development work.
\end{rDated}
\begin{rDated}{Jan. 2022 -- \\ Feb. 2023}{Research Intern}{Seoul, KR}{KIST AI\textperiodcentered Robotics Institute \textbar\ Vision Intelligence Group}
  \item Began surveillance video research through urban public safety projects and national research programs.
\end{rDated}
\begin{rDated}{Oct. 2021 -- \\ Dec. 2021}{Researcher}{Sejong, KR}{The Korea Transport Institute}
  \item Contributed to autonomous bus development through SLAM, control center systems, policy research, and coordination.
\end{rDated}
\end{rSection}

\begin{rSection}{Selected Publications}
\begin{enumerate}[leftmargin=1.7em,label={[\arabic*]},itemsep=7pt,parsep=0pt,topsep=2pt]
  \publication
    {\href{https://skill-lens.net/}{Look Closer: Action-Guided Dense Visual Dynamics for Proficiency Estimation.}}
    {\textbf{Ahn D.}, Kim H., Han J.-H., Ham B., Choi H., Kim I.-J., Kim H.}
    {CVPRW 2026 \textbar\ Multimodal Human Motion Analysis \textbar\ Accepted; Pitch and Poster Presentation}
  \publication
    {\href{https://deo-ahn.github.io/buster/}{BUSTER: Adaptive Sampling for VLM-Guided Unsupervised Video Anomaly Detection.}}
    {\textbf{Ahn D.}, Choi Y., Han J.-H., Kim I.-J., Kim H.}
    {British Machine Vision Conference (BMVC), 2026 \textbar\ Accepted}
  \publication
    {Learning Robust Representations for Few-Shot Action Recognition with Frame-Level Ambiguities.}
    {\textbf{Ahn D.}, Jo Y., Lee U.-S., Park H., Han J.-H., Kim H.}
    {IEEE Access, 2026 \textbar\ Accepted}
  \publication
    {\href{https://ieeexplore.ieee.org/document/11217333}{Where and What: Contextual Dynamics-Aware Anomaly Detection in Surveillance Videos.}}
    {\textbf{Ahn D.}, Jo Y., Kim D., Nam G. P., Han J.-H., Kim H.}
    {IEEE Transactions on Image Processing, 2025 \textbar\ Published\\
    JIF 15.3 \textbar\ Top 2.17\% in E.E. Eng. (2025 JCR)}
  \publication
    {A Foundational Research Framework for Real-World Abandoned Object Detection: Train-Free Baseline and a Standardized Benchmark.}
    {Kim D., \textbf{Ahn D.}, Jo Y., Park H., Lee S., Kim H.}
    {Expert Systems with Applications, 2025 \textbar\ 130658 \textbar\ Published}
  \publication
    {A Passive Tracking System Based on Geometric Constraints in Adaptive Wireless Sensor Networks.}
    {Zhou B., \textbf{Ahn D.}, Lee J., Sun C., Ahmed S., Kim Y.}
    {Sensors, 18(10), 3276, 2018 \textbar\ Published}
  \publication
    {2D Indoor Map Building Scheme Using Ultrasonic Module.}
    {\textbf{Ahn D.}, Kim N. M., Park J. H., Kim Y. O.}
    {J. Korean Inst. Comm. and Info. Sci., 41(8), 986--994, 2016 \textbar\ Published}
\end{enumerate}

\end{rSection}

\newpage
\begin{rSection}{Research Grants \& Projects}
\begin{rDated}{\textasciitilde\$24M}{Real-Time Intelligent Surveillance for Public Safety}{}{National R\&D \textbar\ Five year program}
  \item MSIT \& MOTIE; with Yonsei University and the Korean National Police Agency.
  \item Video anomaly detection that uses context across city scale CCTV control centers with about
  1,000 cameras, validated in Anyang City.
\end{rDated}
\begin{rDated}{\textasciitilde\$9.5M}{Mobile Platform for Active Safety Services}{}{Korean National Police Agency \& MOTIE \textbar\ Five year program}
  \item Government research program using mobile platforms for proactive public safety services.
\end{rDated}
\begin{rDated}{New \\ initiative}{Anomaly Detection in Correctional Facilities}{}{KIST AI\textperiodcentered Robotics}
  \item Detecting inmate anomalous behavior in custodial environments.
\end{rDated}
\end{rSection}

\begin{rSection}{Awards \& Honors}
\begin{rDated}{Dec. 2025}{AIR Outstanding Research Award}{}{KIST AI\textperiodcentered Robotics Institute}
  \item In recognition of outstanding research achievements during 2025.
\end{rDated}
\begin{rDated}{Feb. 2025}{Excellent Paper Award, Graduate School}{}{Korea University}
  \item For outstanding research performance and publication of an excellent paper during doctoral studies.
\end{rDated}
\begin{rDated}{Nov. 2024}{Outstanding Oral Presentation Award}{}{KIST Academia and Research Convergence Conference}
  \item For the presentation ``Multi-Aspect Contextual Dynamics for Anomaly Detection in Surveillance Videos.''
\end{rDated}
\begin{rDated}{Jul. 2021}{Division Commander's Commendation}{}{Capital Mechanized Infantry Division}
  \item For a combat development proposal on reconnaissance drone AI improvement and acoustic detection using AI and big data.
\end{rDated}
\end{rSection}

\begin{rSection}{Teaching}
\begin{rDated}{Sep. -- \\ Dec. 2018}{Senior Capstone Project}{}{Teaching Assistant \textbar\ Kwangwoon University \textbar\ Instructor: Prof. Youngok Kim}
  \item Assisted in managing and supporting the Capstone Design Project as a TA for 50+ students.
\end{rDated}
\begin{rDated}{Mar. -- \\ Jun. 2017}{Computer Architecture}{}{Teaching Assistant \textbar\ Kwangwoon University \textbar\ Instructor: Prof. Youngok Kim}
  \item Graded assignments and provided feedback for 30+ students with Q\&A sessions.
\end{rDated}
\end{rSection}

\begin{rSection}{Technical Skills}
\begin{tabular}{@{}p{1.6in}@{}p{\dimexpr\linewidth-1.6in\relax}@{}}
\textbf{Programming} & Python, C, Java, Git, Linux \\
\textbf{Machine Learning} & PyTorch, TensorFlow, Keras, NumPy, Pandas, Matplotlib \\
\textbf{Sensors} & Ultrasonic, Beacon, RF, Wi-Fi, 3-axis Accelerometer \\
\textbf{Embedded \& IoT} & Raspberry Pi, Arduino \\
\end{tabular}
\end{rSection}

\vfill
\hfill{\small Updated October 2026}
\end{document}
